%% RESUMO NA LINGUA VERNACULA
\begin{resumo}
Esta dissertação dummy organiza o esqueleto de uma pesquisa sobre o uso de inteligência artificial na triagem de manifestações dirigidas às ouvidorias públicas. O objetivo geral é analisar, de forma exploratória, como técnicas de processamento de linguagem natural podem apoiar a classificação e o encaminhamento de demandas, preservando a supervisão humana e a rastreabilidade das decisões. O referencial teórico reúne três estudos recentes sobre classificação de petições, roteamento de denúncias e orquestração de agentes com humano no circuito. O procedimento previsto é uma revisão bibliográfica estruturada em bases abertas, com leitura dos resumos e comparação do contexto institucional e da tarefa tratada. Por se tratar de um esqueleto, o trabalho não apresenta resultados próprios e indica como próximos passos a ampliação da revisão e a definição de critérios de seleção.
\end{resumo}
